\documentclass[spanish,letterpaper,oneside]{article}

\def\documenttitle {Entrada y salida de datos}
\def\documentsubtitle {Actividad 2: fundamentos de programación en C++}
\def\documentsubject {Análisis, algoritmos y codificación para Dev-C++}
\def\documentauthor {Emilio Abraham Castillo Urdiales}
\def\coursename {Aplicación práctica de la programación}
\def\coursecode {Grupo 011 A2026}
\def\universityname {Universidad Autónoma de Nuevo León}
\def\universityfaculty {Facultad de Ingeniería Mecánica y Eléctrica}
\def\universitydepartment {Ingeniería en Administración y Sistemas}
\def\universitydepartmentimage {UANL/assets-uanl/logo-uanl.png}
\def\universitydepartmentimagecfg {width=3.2cm}
\def\universitylocation {San Nicolás de los Garza, Nuevo León, México}
\def\authortable {
  \begin{tabular}{ll}
    \textbf{Alumno:} & Emilio Abraham Castillo Urdiales \\
    Matrícula: & 2194561 \\
    Grupo: & 011 A2026 \\
    Programa: & Ingeniería en Administración y Sistemas \\
    Materia: & \coursename \\
    Docente: & Karla Patricia Uribe Sierra \\
    Periodo: & Agosto--diciembre de 2026 \\
    & \\
    \multicolumn{2}{l}{Fecha de realización: \today} \\
    \multicolumn{2}{l}{\universitylocation}
  \end{tabular}
}

\input{template}
\newcommand{\codepath}{UANL/ingenieria-en-administracion-y-sistemas/metodologia-de-la-programacion/programas-actividad-2}
\definecolor{consolebg}{RGB}{28,30,34}
\definecolor{consolefg}{RGB}{232,235,239}
\lstdefinestyle{console}{
  style=plaintext,
  backgroundcolor=\color{consolebg},
  basicstyle=\ttfamily\footnotesize\color{consolefg},
  rulecolor=\color{black},
  frame=single,
  framerule=0.6pt,
  breaklines=true,
  showstringspaces=false
}
\newcommand{\consola}[1]{%
  \lstinputlisting[style=console,caption={Salida de consola.}]{\codepath/#1.out}%
}
\newcommand{\programa}[5]{%
  \section{#1}
  \subsection{Instrucción}
  #2
  \subsection{Análisis del problema}
  #3
  \subsection{Algoritmo}
  \begin{enumerate}
    #4
  \end{enumerate}
  \subsection{Codificación del programa}
  \importsourcecode{cpp}{\codepath/#5}{Código C++ compatible con Dev-C++.}
  \subsection{Salida de consola}
  \consola{#5}
}

\begin{document}
\templatePortrait
\templatePagecfg
\setlength{\footskip}{28pt}
\fancyfoot[L]{}
\fancyfoot[C]{\footnotesize\thepage}
\fancyfoot[R]{}
\onehalfspacing

\begin{abstractd}
Se resuelven diecinueve problemas de entrada, proceso y salida en C++. Cada solución identifica datos, operación matemática, validaciones y resultado. Los archivos son independientes y compatibles con el compilador GNU C++ de Dev-C++. Los tipos \texttt{double} conservan la precisión decimal; las constantes identifican conversiones y magnitudes físicas; y los mensajes orientan la captura de datos.
\end{abstractd}

\templateIndex
\templateFinalcfg

\section{Introducción}
El modelo entrada--proceso--salida organiza cada solución. La entrada reúne los valores del usuario, el proceso aplica fórmulas o conversiones y la salida comunica el resultado. El análisis previo determina datos, unidades, ecuaciones y condiciones de validez para problemas geométricos, físicos, de conversión y de administración.

\programa{Programa 1. Área de un cuadrado}
{Calcular e imprimir el área de un cuadrado a partir de la longitud de su lado.}
{La entrada es el lado $l$, que debe ser positivo. El proceso aplica $A=l\times l$ y la salida es el área en unidades cuadradas.}
{\item Solicitar el lado. \item Verificar que sea mayor que cero. \item Calcular $A=l\times l$. \item Mostrar el área con dos decimales.}
{01_area_cuadrado.cpp}

\programa{Programa 2. Área de un cuadrado con \texttt{pow}}
{Calcular el área de un cuadrado mediante la función de potencia.}
{Se recibe el lado $l$ y se eleva al cuadrado con \texttt{pow(lado, 2)}. El uso de \texttt{<cmath>} permite practicar funciones matemáticas de la biblioteca estándar.}
{\item Leer el lado. \item Validar que sea positivo. \item Calcular $A=l^2$ con \texttt{pow}. \item Imprimir el resultado.}
{02_area_cuadrado_pow.cpp}

\programa{Programa 3. Área de un círculo}
{Calcular e imprimir el área de un círculo.}
{La entrada es el radio $r$. Se obtiene una aproximación portable de $\pi$ con $\operatorname{acos}(-1)$ y se aplica $A=\pi r^2$. Un radio no positivo se considera inválido.}
{\item Solicitar el radio. \item Validar $r>0$. \item Calcular $A=\pi r^2$. \item Mostrar el área.}
{03_area_circulo.cpp}

\programa{Programa 4. Área de un triángulo}
{Calcular el área de un triángulo a partir de su base y altura.}
{Las entradas son base $b$ y altura $h$, ambas positivas. El proceso es $A=(b\times h)/2$.}
{\item Leer base y altura. \item Comprobar que sean positivas. \item Calcular $A=bh/2$. \item Imprimir el área.}
{04_area_triangulo.cpp}

\programa{Programa 5. Hipotenusa}
{Calcular la hipotenusa de un triángulo rectángulo.}
{Los datos son los catetos $a$ y $b$. Por el teorema de Pitágoras, la hipotenusa es $c=\sqrt{a^2+b^2}$.}
{\item Leer ambos catetos. \item Validar que sean positivos. \item Calcular $c=\sqrt{a^2+b^2}$. \item Mostrar la hipotenusa.}
{05_hipotenusa.cpp}

\programa{Programa 6. Nombre y semestre}
{Ingresar el nombre del estudiante y el semestre que cursará, e imprimir ambos datos.}
{La entrada combina una cadena de texto y un entero. \texttt{getline} conserva los espacios del nombre completo. Se rechaza un nombre vacío o un semestre no positivo.}
{\item Leer el nombre completo. \item Leer el semestre. \item Validar los datos. \item Imprimir nombre y semestre.}
{06_nombre_semestre.cpp}

\programa{Programa 7. Fahrenheit a centígrados}
{Convertir una temperatura de grados Fahrenheit a grados centígrados.}
{La entrada es $F$ y el proceso utiliza $C=(F-32)\times 5/9$. Se emplean constantes decimales para impedir una división entera.}
{\item Leer la temperatura Fahrenheit. \item Calcular $C=(F-32)5/9$. \item Mostrar Fahrenheit y centígrados.}
{07_fahrenheit_centigrados.cpp}

\programa{Programa 8. Pies a metros}
{Convertir una longitud expresada en pies a metros.}
{Se usa la equivalencia exacta indicada: un pie equivale a $0.3048$ metros. Por tanto, $m=p\times0.3048$.}
{\item Leer la longitud en pies. \item Verificar que no sea negativa. \item Multiplicar por $0.3048$. \item Imprimir los metros.}
{08_pies_metros.cpp}

\programa{Programa 9. Libras a kilogramos}
{Convertir un peso expresado en libras a kilogramos.}
{La entrada son libras y se aplica $kg=lb\times0.45359237$. La salida conserva cuatro decimales para mostrar adecuadamente la conversión.}
{\item Leer las libras. \item Validar que no sean negativas. \item Multiplicar por $0.45359237$. \item Mostrar los kilogramos.}
{09_libras_kilogramos.cpp}

\programa{Programa 10. Acres a hectáreas}
{Introducir la extensión de un terreno en acres y calcularla en hectáreas, considerando $1$ acre $=4047\,m^2$ y $1$ hectárea $=10000\,m^2$.}
{Primero se convierten acres a metros cuadrados y después se dividen entre los metros cuadrados de una hectárea: $ha=acres\times4047/10000$.}
{\item Leer acres. \item Validar que no sean negativos. \item Convertir a metros cuadrados. \item Dividir entre $10000$. \item Imprimir hectáreas.}
{10_acres_hectareas.cpp}

\programa{Programa 11. Fuerza}
{Calcular la fuerza aplicada a un cuerpo.}
{De acuerdo con la segunda ley de Newton, $F=m\times a$. Se ingresan masa en kilogramos y aceleración en $m/s^2$; el resultado se expresa en newtons.}
{\item Leer masa y aceleración. \item Validar que la masa no sea negativa. \item Calcular $F=ma$. \item Mostrar la fuerza en newtons.}
{11_fuerza.cpp}

\programa{Programa 12. Velocidad final de impacto}
{Calcular la velocidad final de impacto de un cuerpo en caída libre desde el reposo.}
{Sin resistencia del aire y con velocidad inicial cero, se aplica $v=\sqrt{2gh}$, con $g=9.81\,m/s^2$ y altura $h$ no negativa.}
{\item Leer la altura. \item Verificar $h\geq0$. \item Calcular $v=\sqrt{2gh}$. \item Imprimir la velocidad en $m/s$.}
{12_velocidad_impacto.cpp}

\programa{Programa 13. Energía potencial}
{Calcular la energía potencial de un cuerpo en caída libre.}
{Las entradas son masa $m$ y altura $h$. Con $g=9.81\,m/s^2$, la energía potencial gravitatoria es $E_p=mgh$ y se expresa en joules.}
{\item Leer masa y altura. \item Validar que no sean negativas. \item Calcular $E_p=mgh$. \item Mostrar joules.}
{13_energia_potencial.cpp}

\programa{Programa 14. Energía cinética}
{Calcular la energía cinética de un cuerpo en movimiento.}
{Se ingresan masa $m$ y velocidad $v$. El cálculo corresponde a $E_c=mv^2/2$. La velocidad puede tener signo, pero al elevarla al cuadrado la energía permanece no negativa.}
{\item Leer masa y velocidad. \item Validar la masa. \item Calcular $E_c=mv^2/2$. \item Mostrar joules.}
{14_energia_cinetica.cpp}

\programa{Programa 15. Hipotenusa dados los catetos}
{Dados los catetos de un triángulo rectángulo, calcular e imprimir su hipotenusa.}
{Esta variante utiliza \texttt{hypot}, función estándar que calcula directamente $\sqrt{a^2+b^2}$ y mejora la claridad del código.}
{\item Leer los catetos. \item Comprobar que sean positivos. \item Evaluar \texttt{hypot(a,b)}. \item Imprimir el resultado.}
{15_hipotenusa_catetos.cpp}

\programa{Programa 16. Volumen de una esfera}
{Calcular e imprimir el volumen de una esfera.}
{La entrada es el radio $r>0$. El volumen se obtiene con $V=(4/3)\pi r^3$. Se escribe \texttt{4.0/3.0} para conservar la parte decimal.}
{\item Leer el radio. \item Validar que sea positivo. \item Calcular $V=(4/3)\pi r^3$. \item Mostrar unidades cúbicas.}
{16_volumen_esfera.cpp}

\programa{Programa 17. Recibo de la compañía de luz}
{Generar el recibo de un cliente con su nombre, consumo en kilowatts y tarifa por kilowatt; el recibo debe mostrar el nombre y el pago.}
{El pago se calcula como $P=consumo\times tarifa$. Se usa \texttt{getline} para aceptar nombres completos y se impiden consumo o tarifa negativos.}
{\item Leer el nombre del cliente. \item Leer consumo y tarifa. \item Validar las entradas. \item Calcular el pago. \item Imprimir nombre y total con dos decimales.}
{17_recibo_luz.cpp}

\programa{Programa 18. Conversión conjunta de pies y libras}
{Introducir longitud en pies y peso en libras, e imprimir metros y kilogramos.}
{Se resuelven dos conversiones independientes: $m=pies\times0.3048$ y $kg=libras\times0.45359237$.}
{\item Leer pies y libras. \item Verificar que no sean negativos. \item Convertir pies a metros. \item Convertir libras a kilogramos. \item Mostrar ambos resultados.}
{18_pies_libras.cpp}

\programa{Programa 19. Área mediante la fórmula de Herón}
{Leer los tres lados de un triángulo y calcular su área con la fórmula de Herón.}
{Antes del cálculo se valida la desigualdad triangular. Si los lados son válidos, el semiperímetro es $s=(a+b+c)/2$ y el área $A=\sqrt{s(s-a)(s-b)(s-c)}$.}
{\item Leer $a$, $b$ y $c$. \item Verificar positividad y desigualdad triangular. \item Calcular $s=(a+b+c)/2$. \item Calcular la raíz de Herón. \item Imprimir el área.}
{19_area_heron.cpp}

\section{Pseudocódigo completo}
Los diecinueve algoritmos reproducen la entrada, validación, cálculo y salida de cada programa en sintaxis PSeInt.
\importsourcecode{plaintext}{\codepath/pseudocodigos-actividad-2.psc}{Pseudocódigo PSeInt de los 19 problemas.}

% El script automatizado de compilación y pruebas se conserva como herramienta
% interna del proyecto, pero no forma parte del contenido académico del reporte.
% La sección de pruebas recomendadas también se omite.

\section{Resumen}
\begin{framed}
\textbf{Entrada, proceso y salida.} La programación expresa procedimientos mediante instrucciones precisas y ordenadas. Cada solución inicia con el análisis de las entradas, continúa con una fórmula o transformación y termina con una salida comprensible. El tipo \texttt{double} conserva valores decimales y las constantes con nombre identifican factores como la gravedad y las equivalencias de unidades.

\medskip
\textbf{Bibliotecas y validación.} \texttt{iostream} administra la entrada y salida; \texttt{iomanip} controla los decimales; \texttt{cmath} proporciona potencias y raíces; y \texttt{string} almacena nombres completos. Las validaciones impiden cálculos con lados negativos, radios nulos o longitudes que no forman un triángulo.

\medskip
\textbf{Aplicación.} Formular el algoritmo antes del código reduce errores y permite verificar la unidad y el propósito de cada variable. Los programas se compilan por separado en Dev-C++ y sus salidas de consola confirman los resultados de los casos evaluados.
\end{framed}

\section{Conclusiones}
Los diecinueve programas muestran que el modelo entrada--proceso--salida permite resolver problemas diversos con una estructura común. El análisis previo ayudó a seleccionar fórmulas, tipos de datos y validaciones. El uso de funciones de la biblioteca estándar hizo los cálculos más claros y portables. Como mejora posterior pueden añadirse ciclos para repetir cada operación y controles para detectar entradas que no sean numéricas; sin embargo, las versiones entregadas satisfacen el objetivo de practicar instrucciones básicas y pueden compilarse individualmente en Dev-C++.

\end{document}
